\section{事故機構・設計運用と個別責任の区別}

\subsection{平成30年度 問14}
\label{q:h30-reg-apt-14}

\noindent\textbf{原典：} sources/H30/aptitude.pdf、PDF物理14頁。

\textbf{公式正答：} ALL（all\_awarded）。

\par\textbf{独立解答と公式正答に相違あり：} 独立解答 2|5。初回2|5に対し公式ALL。2の誤った有罪は最高裁2017年決定で照合、5の不適格は原典公式注記で確認。1・3・4の個別判決・量刑比較は一次裁判資料未確認。全肢真とせず部分校正、Astra難問査読へ。

\paragraph{問題の要旨} 5事故の刑事責任を判決と比較する設問と、笹子の不起訴を判決扱いした不適切出題注記を検討する。

\Needspace{4\baselineskip}\paragraph{確認した条件・記述（独自要約）}\mbox{}\par

\noindent\textbf{方向：} 実際の刑事判決と異なる事例を選ぶ。\par

\noindent\textbf{原典注記：} 選択肢5は不起訴処分で判決ではなく、試験日現在検察審査会への申立てがあり不適切出題。\par

\noindent\textbf{時点：} 平成30年の試験日。後年の経過から逆算しない。\par

\noindent\textbf{OCR補正：} 選択肢1はシンドラー社元社員。\par

\paragraph{問われる核心} 事故の技術的原因と個別の刑事責任・採点措置の区別

\paragraph{必要知識} 事故の技術的原因と個別の刑事責任・採点措置の区別。初回独立候補2|5を保持する。2の有罪は、歴代3社長を対象とした最高裁の2017年決定と異なる。5は、公式問題の注記が示すとおり不起訴処分を刑事裁判の判決と同列に置いている。公式正答は不適切出題による全員正解（ALL）であり、初回候補をALLへ上書きしない。1・3・4の個別判決の一次資料は未確認で、全肢を確定済みにしない。全員正解は採点上の措置であり、5事例の記述が全て正しいという意味ではない。

\paragraph{考え方・正答を選ぶ理由（査読待ち初稿）} 初回独立候補2|5を保持する。2の有罪は、歴代3社長を対象とした最高裁の2017年決定と異なる。5は、公式問題の注記が示すとおり不起訴処分を刑事裁判の判決と同列に置いている。公式正答は不適切出題による全員正解（ALL）であり、初回候補をALLへ上書きしない。1・3・4の個別判決の一次資料は未確認で、全肢を確定済みにしない。全員正解は採点上の措置であり、5事例の記述が全て正しいという意味ではない。

\Needspace{4\baselineskip}\paragraph{選択肢の検討（査読待ち初稿）}\mbox{}\par

\noindent\textbf{1：判決一次資料要確認} 2006年エレベーター事故で、シンドラー社元社員の刑事責任はないとされた。。原画像の主語は元社員であり、OCRの社員を補正する。個別被告人の判決と確定段階を示す一次の裁判資料は未確認であり、会社・全保守関係者の責任へ一般化しない。\par

\noindent\textbf{2：判決と異なる} 2005年福知山線事故で、歴代3社長が有罪となった。。2017年6月12日最高裁第二小法廷決定、平成27年（あ）741号は、歴代3社長への無罪の第一審を是認した原判断を相当とし、各上告を棄却した。有罪という結末が異なる。\par

\noindent\textbf{3：判決一次資料要確認} 2004年回転ドア事故で、メーカー営業開発部長の刑事責任が森ビル関係者より重いとされた。。事故防止上の製造者・管理者の役割と、個別被告人への量刑比較は別問題である。刑の重さと理由を一次の判決文で照合できておらず、適切と確定しない。\par

\noindent\textbf{4：判決一次資料要確認} 2000年低脂肪乳食中毒事故で、工場関係者が有罪となった。。厚労省資料は停電時の菌増殖・エンテロトキシンA産生と脱脂粉乳の因果を示すが、工場関係者の有罪判決そのものではない。刑事事件の一次記録は追加確認を要する。\par

\noindent\textbf{5：公式注記による不適格} 2012年笹子トンネル事故で、社長らは予見できないとして刑事責任なしとする判決を受けた。。原典に付された公式説明は、不起訴処分で刑事裁判に当たらず、試験日現在は検察審査会への審査申立てもあったため不適格とする。不起訴を無罪判決に読み替えられない。\par

\paragraph{再利用する考え方} 対象、条件、時点を先に確認する。事故の技術的原因と個別の刑事責任・採点措置の区別を具体的な事例へ適用し、技術的原因・法的要件・倫理的判断のどれを問うか区別する。

\paragraph{共通理論} \hyperref[theory:safety:responsibility]{第\ref*{theory:safety:responsibility}節}, \hyperref[theory:safety:cases]{第\ref*{theory:safety:cases}節}

\par 参考：\cite{court_jr2017,jtsb_fukuchiyama2007,mhlw_milk2002,mlit_sasago2013}。

\paragraph{関連問題} \hyperref[q:r05-reg-apt-09]{R05-REG-APT-09}

\par 難易度 3（校準初稿）、優先度未評価。

\subsection{令和5年度 問9}
\label{q:r05-reg-apt-09}

\noindent\textbf{原典：} sources/R05/aptitude.pdf、PDF物理9頁。

\textbf{公式正答：} 4（single）。

\paragraph{問題の要旨} 食中毒、回転ドア、鉄道、エレベーター、トンネル天井板の事故原因を比較する。

\Needspace{4\baselineskip}\paragraph{確認した条件・記述（独自要約）}\mbox{}\par

\noindent\textbf{方向：} 事実と異なるもの。\par

\noindent\textbf{事故年：} 2000、2004、2005、2006、2012。\par

\noindent\textbf{条件：} 鉄道のATSは当該急曲線の速度超過防止。エレベーターの運動方向を実画像と政府事故調査で確認。\par

\paragraph{問われる核心} 事故機構の再現と設計・運用・組織の因果の区別

\paragraph{必要知識} 事故機構の再現と設計・運用・組織の因果の区別。4は事故の機構を逆にしている。2006年の事象は、かごが戸開状態で不意に上昇した際の挟まれであり、ケーブル破断・落下ではない。1〜3・5は、直接の物理的事象だけでなく、その前提を維持できなかった設計・運用・組織上の条件も述べる。鉄道のATSの機能と、組織背景の推定の強さを区別して4を選ぶ。

\paragraph{考え方・正答を選ぶ理由（査読待ち初稿）} 4は事故の機構を逆にしている。2006年の事象は、かごが戸開状態で不意に上昇した際の挟まれであり、ケーブル破断・落下ではない。1〜3・5は、直接の物理的事象だけでなく、その前提を維持できなかった設計・運用・組織上の条件も述べる。鉄道のATSの機能と、組織背景の推定の強さを区別して4を選ぶ。

\Needspace{4\baselineskip}\paragraph{選択肢の検討（査読待ち初稿）}\mbox{}\par

\noindent\textbf{1：適切（事故対応の説明）} 2000年食中毒は停電後の不適切な処置と情報開示・回収の遅れ等が関係した。。厚労省は停電時の菌増殖・毒素産生を説明し、当事企業も汚染原料と対応の問題を振り返る。個別の経営者の刑事責任を認定する記述ではない。\par

\noindent\textbf{2：適切（安全の教訓）} 2004年回転ドアは検知・停止の設計と運用、軽量設計の軽視等に問題があった。。検知・非常停止だけに安全を委ね、重量・慣性と使用条件を軽視すると停止前の挟まれ危険が残る。大学の技術者倫理事例集は軽量設計からの変更を扱う。事故の一般的教訓と個別被告人の判決を分ける。\par

\noindent\textbf{3：適切（機能・推定の限定に注意）} 2005年福知山線は曲線での速度超過・ブレーキの遅れと安全確保の組織的問題が関係した。。事故調査報告の直接原因は制限速度を大幅に超えた進入で、運転士管理方法の関与は可能性として記す。原典のATS未設置は当該曲線の速度超過防止機能を指し、列車の全ATS不搭載を意味しない。\par

\noindent\textbf{4：事実と異なる} 2006年エレベーター事故はケーブルが破断してかごが落下した。。公式の事故調査は、戸が開いた状態でかごが上昇し利用者が挟まれた事象を扱う。ケーブル破断による落下とした故障機構・運動方向が事実と異なる。\par

\noindent\textbf{5：適切（複合的要因）} 2012年笹子は天井板の接着系アンカー等の脱落、設計・施工・維持管理が関係した。。国交省委員会は接着系アンカーの引抜き等と、設計・施工・点検の課題を整理する。単純な材料の経年劣化だけに原因を限定しない。\par

\paragraph{再利用する考え方} 対象、条件、時点を先に確認する。事故機構の再現と設計・運用・組織の因果の区別を具体的な事例へ適用し、技術的原因・法的要件・倫理的判断のどれを問うか区別する。

\paragraph{共通理論} \hyperref[theory:safety:cases]{第\ref*{theory:safety:cases}節}, \hyperref[theory:safety:barriers]{第\ref*{theory:safety:barriers}節}

\par 参考：\cite{mhlw_milk2002,milk_incident_company,muroran_ethics_cases,jtsb_fukuchiyama2007,caa_elevator2006,mlit_sasago2013}。

\paragraph{関連問題} \hyperref[q:h30-reg-apt-14]{H30-REG-APT-14}

\par 難易度 2（校準初稿）、優先度未評価。

\section{信頼性と故障}

\subsection{平成28年度 問5}
\label{q:h28-reg-apt-05}

\noindent\textbf{原典：} sources/H28/aptitude.pdf、PDF物理5頁。

\textbf{公式正答：} 4（single）。

\par\textit{原典本文の校正と独自解説は未作成。分類のみの登録項目。}

\subsection{令和元年度 問3}
\label{q:r01-reg-apt-03}

\noindent\textbf{原典：} sources/R01/aptitude.pdf、PDF物理4頁。

\textbf{公式正答：} 4（single）。

\par\textit{原典本文の校正と独自解説は未作成。分類のみの登録項目。}

\subsection{令和4年度 問8}
\label{q:r04-reg-apt-08}

\noindent\textbf{原典：} sources/R04/aptitude.pdf、PDF物理7頁。

\textbf{公式正答：} 2（single）。

\par\textit{原典本文の校正と独自解説は未作成。分類のみの登録項目。}

\section{最適化・意思決定}

\subsection{平成26年度 問9}
\label{q:h26-reg-apt-09}

\noindent\textbf{原典：} sources/H26/aptitude.pdf、PDF物理9頁。

\textbf{公式正答：} 4（single）。

\par\textit{原典本文の校正と独自解説は未作成。分類のみの登録項目。}

\subsection{平成27年度 問9}
\label{q:h27-reg-apt-09}

\noindent\textbf{原典：} sources/H27/aptitude.pdf、PDF物理9頁。

\textbf{公式正答：} 3（single）。

\par\textit{原典本文の校正と独自解説は未作成。分類のみの登録項目。}

\subsection{平成28年度 問15}
\label{q:h28-reg-apt-15}

\noindent\textbf{原典：} sources/H28/aptitude.pdf、PDF物理15頁。

\textbf{公式正答：} 3（single）。

\par\textit{原典本文の校正と独自解説は未作成。分類のみの登録項目。}

\subsection{平成29年度 問9}
\label{q:h29-reg-apt-09}

\noindent\textbf{原典：} sources/H29/aptitude.pdf、PDF物理9頁。

\textbf{公式正答：} 3（single）。

\par\textit{原典本文の校正と独自解説は未作成。分類のみの登録項目。}

\subsection{令和4年度 問7}
\label{q:r04-reg-apt-07}

\noindent\textbf{原典：} sources/R04/aptitude.pdf、PDF物理6頁。

\textbf{公式正答：} 3（single）。

\par\textit{原典本文の校正と独自解説は未作成。分類のみの登録項目。}

\subsection{令和5年度 問3}
\label{q:r05-reg-apt-03}

\noindent\textbf{原典：} sources/R05/aptitude.pdf、PDF物理3頁。

\textbf{公式正答：} 2（single）。

\par\textit{原典本文の校正と独自解説は未作成。分類のみの登録項目。}

\subsection{令和7年度 問13}
\label{q:r07-reg-apt-13}

\noindent\textbf{原典：} sources/R07/aptitude.pdf、PDF物理13頁。

\textbf{公式正答：} 5（single）。

\par\textit{原典本文の校正と独自解説は未作成。分類のみの登録項目。}

\section{品質管理・統計}

\subsection{平成23年度 問8}
\label{q:h23-reg-apt-08}

\noindent\textbf{原典：} sources/H23/aptitude.pdf、PDF物理9頁。

\textbf{公式正答：} 5（single）。

\par\textit{原典本文の校正と独自解説は未作成。分類のみの登録項目。}

\subsection{平成25年度 問7}
\label{q:h25-reg-apt-07}

\noindent\textbf{原典：} sources/H25/aptitude.pdf、PDF物理8頁。

\textbf{公式正答：} 4（single）。

\par\textit{原典本文の校正と独自解説は未作成。分類のみの登録項目。}

\subsection{平成30年度 問9}
\label{q:h30-reg-apt-09}

\noindent\textbf{原典：} sources/H30/aptitude.pdf、PDF物理9頁。

\textbf{公式正答：} 3（single）。

\par\textit{原典本文の校正と独自解説は未作成。分類のみの登録項目。}

\subsection{令和元年度 問7}
\label{q:r01-reg-apt-07}

\noindent\textbf{原典：} sources/R01/aptitude.pdf、PDF物理8頁。

\textbf{公式正答：} 5（single）。

\par\textit{原典本文の校正と独自解説は未作成。分類のみの登録項目。}

\section{アルゴリズム・再帰}

\subsection{平成28年度 問7}
\label{q:h28-reg-apt-07}

\noindent\textbf{原典：} sources/H28/aptitude.pdf、PDF物理7頁。

\textbf{公式正答：} 5（single）。

\par\textit{原典本文の校正と独自解説は未作成。分類のみの登録項目。}

\section{材料の応力・変形}

\subsection{令和元年度 問13}
\label{q:r01-reg-apt-13}

\noindent\textbf{原典：} sources/R01/aptitude.pdf、PDF物理13頁。

\textbf{公式正答：} 4（single）。

\par\textit{原典本文の校正と独自解説は未作成。分類のみの登録項目。}

\section{環境法・循環・持続可能性}

\subsection{令和7年度 問15}
\label{q:r07-reg-apt-15}

\noindent\textbf{原典：} sources/R07/aptitude.pdf、PDF物理15頁。

\textbf{公式正答：} 2（single）。

\par\textit{原典本文の校正と独自解説は未作成。分類のみの登録項目。}

\section{公益・説明責任}

\subsection{平成26年度 問11}
\label{q:h26-reg-apt-11}

\noindent\textbf{原典：} sources/H26/aptitude.pdf、PDF物理11頁。

\textbf{公式正答：} 1（single）。

\par\textit{原典本文の校正と独自解説は未作成。分類のみの登録項目。}

\subsection{平成26年度 問15}
\label{q:h26-reg-apt-15}

\noindent\textbf{原典：} sources/H26/aptitude.pdf、PDF物理15頁。

\textbf{公式正答：} 3（single）。

\par\textit{原典本文の校正と独自解説は未作成。分類のみの登録項目。}

\subsection{平成27年度 問4}
\label{q:h27-reg-apt-04}

\noindent\textbf{原典：} sources/H27/aptitude.pdf、PDF物理4頁。

\textbf{公式正答：} 3（single）。

\par\textit{原典本文の校正と独自解説は未作成。分類のみの登録項目。}

\subsection{平成27年度 問5}
\label{q:h27-reg-apt-05}

\noindent\textbf{原典：} sources/H27/aptitude.pdf、PDF物理5頁。

\textbf{公式正答：} 5（single）。

\par\textit{原典本文の校正と独自解説は未作成。分類のみの登録項目。}

\subsection{平成28年度 問3}
\label{q:h28-reg-apt-03}

\noindent\textbf{原典：} sources/H28/aptitude.pdf、PDF物理3頁。

\textbf{公式正答：} 1（single）。

\par\textit{原典本文の校正と独自解説は未作成。分類のみの登録項目。}

\subsection{平成28年度 問9}
\label{q:h28-reg-apt-09}

\noindent\textbf{原典：} sources/H28/aptitude.pdf、PDF物理9頁。

\textbf{公式正答：} 5（single）。

\par\textit{原典本文の校正と独自解説は未作成。分類のみの登録項目。}

\subsection{令和元年度再試験 問5}
\label{q:r01-ret-apt-05}

\noindent\textbf{原典：} sources/R01\_retest/aptitude.pdf、PDF物理5頁。

\textbf{公式正答：} 3（single）。

\par\textit{原典本文の校正と独自解説は未作成。分類のみの登録項目。}

\subsection{令和3年度 問2}
\label{q:r03-reg-apt-02}

\noindent\textbf{原典：} sources/R03/aptitude.pdf、PDF物理2頁。

\textbf{公式正答：} 4（single）。

\par\textit{原典本文の校正と独自解説は未作成。分類のみの登録項目。}

\section{安全・リスク評価}

\subsection{平成23年度 問5}
\label{q:h23-reg-apt-05}

\noindent\textbf{原典：} sources/H23/aptitude.pdf、PDF物理6頁。

\textbf{公式正答：} 4（single）。

\par\textit{原典本文の校正と独自解説は未作成。分類のみの登録項目。}

\subsection{平成23年度 問6}
\label{q:h23-reg-apt-06}

\noindent\textbf{原典：} sources/H23/aptitude.pdf、PDF物理7頁。

\textbf{公式正答：} 4（single）。

\par\textit{原典本文の校正と独自解説は未作成。分類のみの登録項目。}

\subsection{平成24年度 問10}
\label{q:h24-reg-apt-10}

\noindent\textbf{原典：} sources/H24/aptitude.pdf、PDF物理9頁。

\textbf{公式正答：} 1（single）。

\par\textit{原典本文の校正と独自解説は未作成。分類のみの登録項目。}

\subsection{平成24年度 問6}
\label{q:h24-reg-apt-06}

\noindent\textbf{原典：} sources/H24/aptitude.pdf、PDF物理5頁。

\textbf{公式正答：} 3（single）。

\par\textit{原典本文の校正と独自解説は未作成。分類のみの登録項目。}

\subsection{平成25年度 問4}
\label{q:h25-reg-apt-04}

\noindent\textbf{原典：} sources/H25/aptitude.pdf、PDF物理5頁。

\textbf{公式正答：} 5（single）。

\par\textit{原典本文の校正と独自解説は未作成。分類のみの登録項目。}

\subsection{平成25年度 問8}
\label{q:h25-reg-apt-08}

\noindent\textbf{原典：} sources/H25/aptitude.pdf、PDF物理9頁。

\textbf{公式正答：} 2（single）。

\par\textit{原典本文の校正と独自解説は未作成。分類のみの登録項目。}

\subsection{平成25年度 問9}
\label{q:h25-reg-apt-09}

\noindent\textbf{原典：} sources/H25/aptitude.pdf、PDF物理10頁。

\textbf{公式正答：} 5（single）。

\par\textit{原典本文の校正と独自解説は未作成。分類のみの登録項目。}

\subsection{平成26年度 問8}
\label{q:h26-reg-apt-08}

\noindent\textbf{原典：} sources/H26/aptitude.pdf、PDF物理8頁。

\textbf{公式正答：} 3（single）。

\par\textit{原典本文の校正と独自解説は未作成。分類のみの登録項目。}

\subsection{平成27年度 問13}
\label{q:h27-reg-apt-13}

\noindent\textbf{原典：} sources/H27/aptitude.pdf、PDF物理13頁。

\textbf{公式正答：} 2（single）。

\par\textit{原典本文の校正と独自解説は未作成。分類のみの登録項目。}

\subsection{平成27年度 問14}
\label{q:h27-reg-apt-14}

\noindent\textbf{原典：} sources/H27/aptitude.pdf、PDF物理14頁。

\textbf{公式正答：} 5（single）。

\par\textit{原典本文の校正と独自解説は未作成。分類のみの登録項目。}

\subsection{平成29年度 問8}
\label{q:h29-reg-apt-08}

\noindent\textbf{原典：} sources/H29/aptitude.pdf、PDF物理8頁。

\textbf{公式正答：} 5（single）。

\par\textit{原典本文の校正と独自解説は未作成。分類のみの登録項目。}

\subsection{平成30年度 問10}
\label{q:h30-reg-apt-10}

\noindent\textbf{原典：} sources/H30/aptitude.pdf、PDF物理10頁。

\textbf{公式正答：} 5（single）。

\par\textit{原典本文の校正と独自解説は未作成。分類のみの登録項目。}

\subsection{平成30年度 問11}
\label{q:h30-reg-apt-11}

\noindent\textbf{原典：} sources/H30/aptitude.pdf、PDF物理11頁。

\textbf{公式正答：} 3（single）。

\par\textit{原典本文の校正と独自解説は未作成。分類のみの登録項目。}

\subsection{令和元年度 問11}
\label{q:r01-reg-apt-11}

\noindent\textbf{原典：} sources/R01/aptitude.pdf、PDF物理12頁。

\textbf{公式正答：} 2（single）。

\par\textit{原典本文の校正と独自解説は未作成。分類のみの登録項目。}

\subsection{令和元年度再試験 問7}
\label{q:r01-ret-apt-07}

\noindent\textbf{原典：} sources/R01\_retest/aptitude.pdf、PDF物理7頁。

\textbf{公式正答：} 3（single）。

\par\textit{原典本文の校正と独自解説は未作成。分類のみの登録項目。}

\subsection{令和2年度 問9}
\label{q:r02-reg-apt-09}

\noindent\textbf{原典：} sources/R02/aptitude.pdf、PDF物理9頁。

\textbf{公式正答：} 3（single）。

\par\textit{原典本文の校正と独自解説は未作成。分類のみの登録項目。}

\subsection{令和3年度 問15}
\label{q:r03-reg-apt-15}

\noindent\textbf{原典：} sources/R03/aptitude.pdf、PDF物理15頁。

\textbf{公式正答：} 3（single）。

\par\textit{原典本文の校正と独自解説は未作成。分類のみの登録項目。}

\subsection{令和4年度 問6}
\label{q:r04-reg-apt-06}

\noindent\textbf{原典：} sources/R04/aptitude.pdf、PDF物理5頁。

\textbf{公式正答：} 5（single）。

\par\textit{原典本文の校正と独自解説は未作成。分類のみの登録項目。}

\subsection{令和5年度 問12}
\label{q:r05-reg-apt-12}

\noindent\textbf{原典：} sources/R05/aptitude.pdf、PDF物理12頁。

\textbf{公式正答：} 4（single）。

\par\textit{原典本文の校正と独自解説は未作成。分類のみの登録項目。}

\subsection{令和6年度 問13}
\label{q:r06-reg-apt-13}

\noindent\textbf{原典：} sources/R06/aptitude.pdf、PDF物理13頁。

\textbf{公式正答：} 5（single）。

\par\textit{原典本文の校正と独自解説は未作成。分類のみの登録項目。}

\subsection{令和6年度 問7}
\label{q:r06-reg-apt-07}

\noindent\textbf{原典：} sources/R06/aptitude.pdf、PDF物理7頁。

\textbf{公式正答：} 4（single）。

\par\textit{原典本文の校正と独自解説は未作成。分類のみの登録項目。}

\subsection{令和7年度 問14}
\label{q:r07-reg-apt-14}

\noindent\textbf{原典：} sources/R07/aptitude.pdf、PDF物理14頁。

\textbf{公式正答：} 2（single）。

\par\textit{原典本文の校正と独自解説は未作成。分類のみの登録項目。}

\section{災害リスク情報・避難制度と事前準備}

\subsection{令和元年度再試験 問13}
\label{q:r01-ret-apt-13}

\noindent\textbf{原典：} sources/R01\_retest/aptitude.pdf、PDF物理14頁。

\textbf{公式正答：} 1（single）。

\paragraph{問題の要旨} 水害・土砂災害で地図、警報、避難情報、事前準備、自主判断の5記述を検討する。

\Needspace{4\baselineskip}\paragraph{確認した条件・記述（独自要約）}\mbox{}\par

\noindent\textbf{方向：} 不適切な記述の数。\par

\noindent\textbf{時点：} 2019年度再試験。旧避難情報名称を現在へ書き換えない。\par

\noindent\textbf{導入部：} AR5の評価した排出シナリオ、1901〜2017日本の降水日数の記述は該当一次資料で確認。すべての仮想シナリオや全地点個別の結果へ拡張しない。\par

\noindent\textbf{ア：} ハザードマップで地域の災害範囲・程度・避難情報を知る。 ○：対象災害や想定条件、経路と場所を事前に理解する。白い範囲が全災害で安全という意味ではない。\par

\noindent\textbf{イ：} 数日前の情報から注意報・警報・特別警報へ危険に応じ発表。 ○：当時の情報体系の概要として適切。ただし急変で常にすべての段階を順次通る保証はない。\par

\noindent\textbf{ウ：} 市町村長の準備・高齢者等避難開始、勧告、指示等で避難。 ○：2019年度の名称として適切。勧告は2021年5月20日に廃止し指示へ一本化。\par

\noindent\textbf{エ：} 備蓄・持出品、安否確認・集合、保険等を事前に準備。 ○：避難、家族との連絡、被害後の生活を異なる対策で支える。\par

\noindent\textbf{オ：} 発令が間に合わなくても危険なら自ら避難判断。 ○：行政情報の受信だけに依存すると急変に間に合わない。自ら得られる危険情報も判断の入力とする。\par

\paragraph{問われる核心} 災害情報を行動に接続する事前準備と制度時点

\paragraph{必要知識} 災害情報を行動に接続する事前準備と制度時点。ア〜オはすべて適切で不適切は0件、候補1。原典の当時の情報名を確認した上で、地図の想定外・突然の危険・警報発表主体を補足する。現在は2021年の避難勧告廃止と、2026年5月29日の警戒レベルを付した新しい防災気象情報を区別して学ぶ。

\paragraph{考え方・正答を選ぶ理由（査読待ち初稿）} ア〜オはすべて適切で不適切は0件、候補1。原典の当時の情報名を確認した上で、地図の想定外・突然の危険・警報発表主体を補足する。現在は2021年の避難勧告廃止と、2026年5月29日の警戒レベルを付した新しい防災気象情報を区別して学ぶ。

\Needspace{4\baselineskip}\paragraph{選択肢の検討（査読待ち初稿）}\mbox{}\par

\noindent\textbf{1：整合する} 0。全5記述が当時の一般的備えとして適切で、不適切0と一致。\par

\noindent\textbf{2：整合しない} 1。不適切な記述を1件とする根拠がない。\par

\noindent\textbf{3：整合しない} 2。不適切な記述を2件とする根拠がない。\par

\noindent\textbf{4：整合しない} 3。不適切な記述を3件とする根拠がない。\par

\noindent\textbf{5：整合しない} 4。不適切な記述を4件とする根拠がない。\par

\paragraph{再利用する考え方} 活動・物質・制度の対象と評価の境界を固定し、用語の入替、全部/一部、法的義務/目標、出題時/現在を区別する。災害情報を行動に接続する事前準備と制度時点を個別の技術活動に結び付ける。

\paragraph{共通理論} \hyperref[theory:safety:preparedness]{第\ref*{theory:safety:preparedness}節}

\par 参考：\cite{env_ar5,env_rain2017,env_evac2021,env_warning2026}。

\par 難易度 2（校準初稿）、優先度未評価。

\section{社会的責任・安全の標準}

\subsection{平成24年度 問11}
\label{q:h24-reg-apt-11}

\noindent\textbf{原典：} sources/H24/aptitude.pdf、PDF物理10頁。

\textbf{公式正答：} 3（single）。

\par\textit{原典本文の校正と独自解説は未作成。分類のみの登録項目。}

\subsection{平成25年度 問5}
\label{q:h25-reg-apt-05}

\noindent\textbf{原典：} sources/H25/aptitude.pdf、PDF物理6頁。

\textbf{公式正答：} 3（single）。

\par\textit{原典本文の校正と独自解説は未作成。分類のみの登録項目。}

\subsection{平成29年度 問12}
\label{q:h29-reg-apt-12}

\noindent\textbf{原典：} sources/H29/aptitude.pdf、PDF物理11頁。

\textbf{公式正答：} 3（single）。

\par\textit{原典本文の校正と独自解説は未作成。分類のみの登録項目。}

\subsection{令和元年度再試験 問12}
\label{q:r01-ret-apt-12}

\noindent\textbf{原典：} sources/R01\_retest/aptitude.pdf、PDF物理13頁。

\textbf{公式正答：} 2（single）。

\par\textit{原典本文の校正と独自解説は未作成。分類のみの登録項目。}

\subsection{令和2年度 問7}
\label{q:r02-reg-apt-07}

\noindent\textbf{原典：} sources/R02/aptitude.pdf、PDF物理7頁。

\textbf{公式正答：} 4（single）。

\par\textit{原典本文の校正と独自解説は未作成。分類のみの登録項目。}

\subsection{令和3年度 問10}
\label{q:r03-reg-apt-10}

\noindent\textbf{原典：} sources/R03/aptitude.pdf、PDF物理10頁。

\textbf{公式正答：} 3（single）。

\par\textit{原典本文の校正と独自解説は未作成。分類のみの登録項目。}

\subsection{令和5年度 問14}
\label{q:r05-reg-apt-14}

\noindent\textbf{原典：} sources/R05/aptitude.pdf、PDF物理14頁。

\textbf{公式正答：} 3（single）。

\par\textit{原典本文の校正と独自解説は未作成。分類のみの登録項目。}

\subsection{令和5年度 問8}
\label{q:r05-reg-apt-08}

\noindent\textbf{原典：} sources/R05/aptitude.pdf、PDF物理8頁。

\textbf{公式正答：} 2（single）。

\par\textit{原典本文の校正と独自解説は未作成。分類のみの登録項目。}

\section{未確定}

\subsection{平成24年度 問7}
\label{q:h24-reg-apt-07}

\noindent\textbf{原典：} sources/H24/aptitude.pdf、PDF物理6頁。

\textbf{公式正答：} 4（single）。

\par\textit{原典本文の校正と独自解説は未作成。分類のみの登録項目。}

\subsection{平成26年度 問13}
\label{q:h26-reg-apt-13}

\noindent\textbf{原典：} sources/H26/aptitude.pdf、PDF物理13頁。

\textbf{公式正答：} 5（single）。

\par\textit{原典本文の校正と独自解説は未作成。分類のみの登録項目。}

\subsection{平成26年度 問7}
\label{q:h26-reg-apt-07}

\noindent\textbf{原典：} sources/H26/aptitude.pdf、PDF物理7頁。

\textbf{公式正答：} 4（single）。

\par\textit{原典本文の校正と独自解説は未作成。分類のみの登録項目。}

\subsection{平成29年度 問5}
\label{q:h29-reg-apt-05}

\noindent\textbf{原典：} sources/H29/aptitude.pdf、PDF物理5頁。

\textbf{公式正答：} 3（single）。

\par\textit{原典本文の校正と独自解説は未作成。分類のみの登録項目。}

\subsection{令和元年度 問12}
\label{q:r01-reg-apt-12}

\noindent\textbf{原典：} sources/R01/aptitude.pdf、PDF物理13頁。

\textbf{公式正答：} 1（single）。

\par\textit{原典本文の校正と独自解説は未作成。分類のみの登録項目。}

\subsection{令和元年度 問8}
\label{q:r01-reg-apt-08}

\noindent\textbf{原典：} sources/R01/aptitude.pdf、PDF物理9頁。

\textbf{公式正答：} 2（single）。

\par\textit{原典本文の校正と独自解説は未作成。分類のみの登録項目。}

\subsection{令和2年度 問13}
\label{q:r02-reg-apt-13}

\noindent\textbf{原典：} sources/R02/aptitude.pdf、PDF物理13頁。

\textbf{公式正答：} 4（single）。

\par\textit{原典本文の校正と独自解説は未作成。分類のみの登録項目。}

\subsection{令和2年度 問6}
\label{q:r02-reg-apt-06}

\noindent\textbf{原典：} sources/R02/aptitude.pdf、PDF物理6頁。

\textbf{公式正答：} 5（single）。

\par\textit{原典本文の校正と独自解説は未作成。分類のみの登録項目。}

\subsection{令和3年度 問8}
\label{q:r03-reg-apt-08}

\noindent\textbf{原典：} sources/R03/aptitude.pdf、PDF物理8頁。

\textbf{公式正答：} 3（single）。

\par\textit{原典本文の校正と独自解説は未作成。分類のみの登録項目。}

\subsection{令和4年度 問11}
\label{q:r04-reg-apt-11}

\noindent\textbf{原典：} sources/R04/aptitude.pdf、PDF物理10頁。

\textbf{公式正答：} 4（single）。

\par\textit{原典本文の校正と独自解説は未作成。分類のみの登録項目。}

\subsection{令和4年度 問5}
\label{q:r04-reg-apt-05}

\noindent\textbf{原典：} sources/R04/aptitude.pdf、PDF物理4頁。

\textbf{公式正答：} 4（single）。

\par\textit{原典本文の校正と独自解説は未作成。分類のみの登録項目。}

\subsection{令和5年度 問6}
\label{q:r05-reg-apt-06}

\noindent\textbf{原典：} sources/R05/aptitude.pdf、PDF物理6頁。

\textbf{公式正答：} 3（single）。

\par\textit{原典本文の校正と独自解説は未作成。分類のみの登録項目。}

\subsection{令和6年度 問14}
\label{q:r06-reg-apt-14}

\noindent\textbf{原典：} sources/R06/aptitude.pdf、PDF物理14頁。

\textbf{公式正答：} 3（single）。

\par\textit{原典本文の校正と独自解説は未作成。分類のみの登録項目。}

\subsection{令和7年度 問12}
\label{q:r07-reg-apt-12}

\noindent\textbf{原典：} sources/R07/aptitude.pdf、PDF物理12頁。

\textbf{公式正答：} 4（single）。

\par\textit{原典本文の校正と独自解説は未作成。分類のみの登録項目。}
